\documentclass[11pt,a4paper]{coverletter}

\senderblock
  {Radheem Bin Razi}
  {Distributed Systems \& Cloud-Native Engineer}
  {Ilmenau, Germany \textbar\ \href{mailto:sheikh.radheem@gmail.com}{sheikh.radheem@gmail.com} \textbar\ \href{https://radheem.github.io/my-cv}{portfolio}}

\begin{document}

%% ===== ENGLISH =====
\recipient{Genki}{Hiring Team}{}

\opening{Dear Hiring Team,}

Genki's pitch — staying healthy wherever you go — lands for me personally: I moved to Germany to start a master's, and reliable cross-border health insurance is exactly the kind of unglamorous-but-essential system I want to help build. Joining as the second developer on a backend team, reading code that isn't fully documented yet, and shaping an event-sourced architecture is the kind of work I've done before and genuinely enjoy.

At Bluefin Exchange I was the design authority for the core services of a high-throughput crypto exchange — a large, complex production backend spanning multiple domains. I ramped quickly on non-obvious business logic, drove event-driven system-design improvements across Go and NodeJS services on Kafka, and set the code-quality bar through design and code reviews. That meant comfort with ambiguity and with deciding before every detail was nailed down — exactly the working style you describe. At Al Hilal Invest I built and documented new backend services in a financial-product domain and tightened the deployment path.

Event-driven architecture is also where my recent work lives. My IRS platform is a set of distributed Go microservices that I migrated off Dapr onto native NATS JetStream and KV with gRPC — event sourcing and append-style messaging at its core. I also build with AI-agentic tools daily: cv-tailor (the app behind this very application) pairs a pure, unit-tested Python ranker with LLM-generated prose, and the IRS platform exposes MCP tooling so LLM agents can operate the system. I care about readable code, meaningful tests, and clean APIs, and I work well async without needing a meeting for everything.

I'm based in Ilmenau, Germany, available full-time now and open to working-student or part-time arrangements. I'm Blue Card-ready, so there's no sponsorship overhead — an employer only issues the contract. A fully remote role with overlap around German business hours fits me perfectly. I'd love to show you what I've got — my GitHub and code samples are ready whenever you are.

\closing{Sincerely,}{Radheem Bin Razi}

\clearpage
\selectlanguage{ngerman}

%% ===== DEUTSCH =====
\recipient{Genki}{Personalabteilung}{}

\opening{Sehr geehrtes Hiring-Team,}

Genkis Anspruch — überall gesund zu bleiben — spricht mich persönlich an: Ich bin nach Deutschland gezogen, um ein Masterstudium zu beginnen, und eine zuverlässige grenzüberschreitende Krankenversicherung ist genau die Art von unspektakulärem, aber unverzichtbarem System, an dessen Aufbau ich mitwirken möchte. Als zweiter Entwickler in ein Backend-Team einzusteigen, noch nicht vollständig dokumentierten Code zu lesen und eine Event-Sourcing-Architektur mitzugestalten, ist genau die Art von Arbeit, die ich bereits gemacht habe und die mir wirklich Freude bereitet.

Bei Bluefin Exchange war ich der Design-Verantwortliche für die Kernservices einer durchsatzstarken Krypto-Börse — ein großes, komplexes Produktions-Backend über mehrere Domänen hinweg. Ich habe mich schnell in nicht offensichtliche Geschäftslogik eingearbeitet, Verbesserungen am ereignisgesteuerten Systemdesign über Go- und NodeJS-Services auf Kafka vorangetrieben und durch Design- und Code-Reviews den Maßstab für die Codequalität gesetzt. Das bedeutete Souveränität im Umgang mit Mehrdeutigkeit und damit, zu entscheiden, bevor jedes Detail feststeht — genau die Arbeitsweise, die Sie beschreiben. Bei Al Hilal Invest habe ich neue Backend-Services in einer Finanzprodukt-Domäne aufgebaut und dokumentiert und den Deployment-Pfad optimiert.

Ereignisgesteuerte Architektur ist auch der Bereich, in dem meine jüngste Arbeit angesiedelt ist. Meine IRS-Plattform besteht aus verteilten Go-Microservices, die ich von Dapr auf natives NATS JetStream und KV mit gRPC migriert habe — mit Event Sourcing und append-orientiertem Messaging im Kern. Außerdem entwickle ich täglich mit KI-agentischen Tools: cv-tailor (die Anwendung hinter genau dieser Bewerbung) verbindet einen reinen, unit-getesteten Python-Ranker mit LLM-generierter Prosa, und die IRS-Plattform stellt MCP-Tooling bereit, sodass LLM-Agenten das System bedienen können. Mir sind lesbarer Code, aussagekräftige Tests und saubere APIs wichtig, und ich arbeite gut asynchron, ohne für alles ein Meeting zu benötigen.

Ich bin in Ilmenau, Deutschland, ansässig, ab sofort in Vollzeit verfügbar und offen für Werkstudenten- oder Teilzeitvereinbarungen. Ich bin Blue-Card-bereit, sodass kein Sponsoring-Aufwand entsteht — ein Arbeitgeber stellt lediglich den Vertrag aus. Eine vollständig remote ausgeübte Rolle mit Überschneidung rund um die deutschen Geschäftszeiten passt perfekt zu mir. Gerne zeige ich Ihnen, was ich mitbringe — mein GitHub und Code-Beispiele stehen bereit, wann immer Sie möchten.

\closing{Mit freundlichen Grüßen,}{Radheem Bin Razi}

\end{document}
